% Source: https://artofproblemsolving.com/community/c3h1978105_mr_p_and_the_clock
% Stable identity: aops-topic-1978105-q01
One possible set (many others exist):
\[
\begin{aligned}
1&=(1+2)\div3, & 2&=1-2+3, & 3&=1^2\times3, & 4&=1^2+3,\\
5&=1\times(2+3), & 6&=1\times2\times3, & 7&=1+2\times3, & 8&=1\times2^3,\\
9&=(1+2)\times3, & 10&=\left\lfloor\sqrt{1\times(2+3)!}\right\rfloor, & 11&=\sqrt{1+(2+3)!}, & 12&=1\times2\times3!.
\end{aligned}
\]
The only hard cases are $10$ and $11$: with $(2+3)!=120$ we get $\sqrt{121}=11$ exactly and $\left\lfloor\sqrt{120}\right\rfloor=10$ because $100\le120<121$.
